\documentclass[12pt]{article}
\usepackage[utf8]{inputenc}
\usepackage[T1]{fontenc}
\usepackage{hyperref}
\title{Sakan Language Resources for K-12 Curriculum Development}
\author{Compiled by Claude Sonnet 5.5 (claude-sonnet-5-5) for Andr\'as Kornai}
\date{}
\begin{document}
\maketitle

\section{Introduction}
Sakan (ISO 639-3 xsk) is not a living language with a school community. Search results indicate that ``Sakan'' refers to Saka, an ancient Eastern Iranian language of the Tarim Basin, attested as Khotanese and Tumshuqese. A 2008 IETF languages mailing-list thread says the code xsk was to be retired as a duplicate, with Sakan treated as a name under Khotanese (kho). The ``Bookkeeping'' family label in the request appears to be a code-management category and not a linguistic family. No K-12 resources exist for it.

\section{Primary Educational Resources}
None found. There is no speaker community and no school use.

\section{Linguistic and Reference Materials}
\begin{itemize}
\item Wikipedia, Saka language: \url{https://en.wikipedia.org/wiki/Saka_language}
\item Glottolog, Saka (saka1298): \url{https://glottolog.org/resource/languoid/id/saka1298.bigmap.html}
\item IETF languages list, ``Khotanese and Sakan'' (2008): \url{https://www.alvestrand.no/pipermail/ietf-languages/2008-September/008375.html}
\end{itemize}

\section{Religious and Community Resources}
None found.

\section{Digital and Audio Resources}
None found.

\section{Curriculum Development Context}
There is no basis for a K-12 curriculum in this language. At most, Saka could appear as a topic in history or linguistics lessons. We recommend that this code be treated as retired or merged, and that any further work be done under Khotanese (kho).

\end{document}